\section{Conclusion}
\label{sec:6}

This work started from a question that an operations manager asks. When a disruption strikes, which reaction achieves the targeted level of resilience, and, in hindsight, was the loss due to the disruption, the network, or the reaction? Answering it required observing the same crises under several management choices, which no real data allow. The approach produces them by simulation, replaying each crisis identically under several policies, measures them with a resilience curve whose parameters are interpretable, and then learns from them a Bayesian network in which the effect of the policy is identifiable.

Over 2,000 crises simulated with Isomorph-Reborn on the ISOMORPH network, the reaction of the operations managers shortens the crisis more than it reduces its depth. The profitability of the levers depends more on the moment at which they are engaged than on their number. When the disruption duration is known, the learned network recommends a lighter policy in about one crisis out of seven, without a notable loss of resilience; reliable information on the duration of an incident is therefore worth seeking.

Several extensions are open. Isomorph-Reborn is fully parametric: the choices and assumptions tied to the network studied here (topology, capacities, delays, inventory policies, profiles, and levers) are not fixed. A user can incorporate the specific characteristics of the network they want to simulate. Two networks react differently to the same disruption and the same decision. Isomorph-Reborn makes it possible to study this effect of topology on varied networks, up to treating network reconfiguration as a resilience lever in its own right. Putting a cost on the levers and on their withdrawal would turn the recommendation into an explicit trade-off between cost and resilience. The critical threshold, which signals a crisis, could be complemented by a survival threshold, below which no recovery is expected any longer and the company may have to cease operations, as the Covid-19 crisis showed. Finally, the recommendations remain to be tested against real crises, even partially documented ones.
